% !TeX root = ../main.tex
\section{Playing techniques}
\label{sec:techniques}

\subsection{What the stage must do}

Mark, per note, how it is played: bend (and how far, in quarter tones) and
release, vibrato, slide into or out of the note, hammer-on and pull-off (a note
sounded by the fretting hand without picking), palm mute, natural and
artificial harmonics, dead (muted) notes. Some of these change the fingering:
a slide or a hammer-on stays on one string, an open string cannot be bent.

\subsection{How techniques are detected}

\begin{keyidea}[: read the pitch contour]
Most techniques are visible in the fine pitch contour of a single note
(Section~\ref{sec:audio}): a bend is a smooth rise and maybe a fall back, vibrato
a periodic wobble, a slide a run of semitone steps, a hammer-on a pitch jump
\emph{without} a new pick attack. Rules on the contour and the attack can
therefore detect them on clean single-note lines.
\end{keyidea}

Chen et al.~\cite{chen2015technique} introduced rule-based detection of bends,
vibrato, slides and hammer-ons on real electric guitar recordings.
TENT~\cite{su2019tent} built technique recognition into note tracking for
monophonic solos, with a 14.7\% higher note F-score than an earlier method
\reported. TART~\cite{gupta2026tart} classifies techniques per note with a small
network (an MLP). GOAT~\cite{loth2025goat} and IDMT-SMT-Guitar~\cite{kehling2014idmt}
provide labels. Almost all evaluation is on clean, monophonic material: distortion
blurs the contour and compression hides the pick attack, so techniques on
amplified rhythm parts are largely untested.

\begin{inapp}[: rules on Basic Pitch's contour \done]
\file{transcribe/techniques.rs}; \code{GVA_TECHNIQUES=0} turns it off.
\begin{itemize}
\item \textbf{Contour}: Basic Pitch's third output, 264 bins at three per
  semitone (MIDI pitch $p$ peaks at bin $3(p-21)+1$, measured). With the CRNN
  engine, Basic Pitch is run again on the guitar audio just for this contour.
\item \textbf{Single-note passages only}; chords are skipped.
\item \textbf{Split bends merged}: the note model often reports a bend as two
  notes; a glide of at least 30\ms{} that does not step through semitones joins
  them back into one bent note.
\item \textbf{Bend and release} inside a note; \textbf{vibrato} at 3.5--9\Hz,
  at least 0.25 semitone deep, at least 2 cycles.
\item \textbf{Slides}: stepped glides, downward glides or glides of more than
  3 semitones, with a gap of up to 120\ms{} between the notes.
\item \textbf{Hammer-on / pull-off}: a pitch jump within 40\ms{} with no
  high-frequency pick burst (\code{PICK_ATTACK}).
\item \textbf{Output}: \code{arrange} keeps techniques per pitch; the tab gets
  \code{bend}, \code{bend_value} (quarter tones) and \code{bend_release} on the
  struck note, \code{slide} and \code{hammer} on the last tied piece, vibrato
  on all pieces. The GP5 export writes bends and shift slides, and the importer
  reads bends back.
\item Thresholds can be swept with \code{GVA_BEND_RISE}, \code{GVA_BEND_START},
  \code{GVA_BEND_GLIDE}.
\end{itemize}
\end{inapp}

\subsection{What we measured}

\begin{itemize}
\item On GuitarSet the rules change nothing measurable (F1 0.83 with or without
  them; 0.76 either way on the crunch render): the 60 excerpts contain only 19
  annotated glides \measured.
\item On real technique labels (IDMT, Guitar-TECHS, GOAT, via
  \code{eval_manifest} with a 5th column in the notes file) \measured:
  vibrato is found on about 20--40\% of labelled vibrato notes, bends on
  35--56\%, slides on about 30\% (often marked as legato instead), and false
  alarms on plain notes are 0--4\%.
\end{itemize}

So the rules are cautious: they rarely invent a technique, and they miss most.

\subsection{Failure modes}
\begin{itemize}
\item No palm mute, harmonics, dead notes or tapping (not implemented).
\item Chords are skipped by design.
\item Under distortion the contour gets noisy and compressed attacks hide the
  pick burst, so bends, vibrato and hammer-ons are missed exactly where they are
  most common \expected.
\end{itemize}

\begin{nextsteps}
Tune the hammer-on attack threshold and the bend rules on GOAT and
Guitar-TECHS; then a learned per-note classifier (task T3 in
Section~\ref{sec:next}) for palm mute, harmonics and dead notes, keeping the
rules for any class where the model is worse.
\end{nextsteps}
